
\section{Related Works}

\subsection{Large Language Models and Retrieval}

Despite demonstrating remarkable reasoning \citep{guo2025deepseek} and coding \citep{guo2024deepseek} capabilities, LLMs \citep{zhao2023survey, team2024gemini, achiam2023gpt} often lack domain-specific knowledge \citep{peng2023study,li2023large} and are prone to hallucinations \citep{zhang2023siren}. To mitigate these limitations, search engines \citep{zhao2024dense} are widely integrated to supply external information.
There are two primary ways to integrate search engines with LLMs: (1) retrieval-augmented generation (RAG) \citep{gao2023retrieval} and (2) treating the search engines as tools \citep{schick2023toolformer}. RAG \citep{lewis2020retrieval, yue2024inference, xiong2025rag} typically follows a round of retrieval and sequential generation pipelines, where a search engine fetches relevant information based on the input query, which is then concatenated with the query and fed into the LLM. However, this could face challenges of retrieving irrelevant information \citep{jin2024long} and failing to provide sufficiently useful context \citep{jiang2023active}.
An alternative approach is search-as-a-tool, where LLMs are prompted or fine-tuned to interact with search engines. IRCoT \citep{trivedi2022interleaving} and ReAct \citep{yao2023react} use prompting to guide iterative reasoning and search engine calls, while Toolformer \citep{schick2023toolformer} leverages supervised fine-tuning to enhance search capabilities. However, such methods rely on high-quality labeled trajectories, which are difficult to obtain at scale. Recent work \citep{guo2025deepseek} suggests that RL can enable LLMs to develop advanced reasoning skills using only outcome rewards, yet its potential in search engine calling scenarios remains under-explored.

% \vspace{-0.1in}
\subsection{Large Language Models and Reinforcement Learning}

Reinforcement learning (RL) \citep{kaelbling1996reinforcement} is a learning paradigm where an agent learns to make sequential decisions by interacting with an environment and receiving feedback in the form of rewards, aiming to maximize cumulative reward over time \citep{sutton1999reinforcement}. RL was introduced to LLM tuning by \cite{ouyang2022training} through RL from human feedback (RLHF) \citep{kaufmann2023survey}. This approach first trains a reward model using human preference data \citep{lambert2024rewardbench}, which then guides RL-based tuning of the policy LLM, typically via Proximal Policy Optimization (PPO). However, PPO involves multiple rounds of LLM optimization, making it challenging to implement.
To simplify RL-based tuning, direct optimization methods such as Direct Preference Optimization (DPO) \citep{rafailov2023direct} and SimPO \citep{meng2024simpo} have been proposed. 
A similar approach is employed in LeRet \citep{hsu2024grounding}, where LLMs are trained to explore diverse queries to enhance the effectiveness of information retrieval.
While these methods offer computational efficiency, they suffer from off-policy issues \citep{pang2024iterative} and do not consistently match the performance of pure RL approaches. Alternative solutions include Group Relative Policy Optimization (GRPO) \citep{shao2024deepseekmath}, which eliminates the need for a critic model by estimating baselines from group scores, and RLOO \citep{ahmadian2024back}, which introduces a simplified REINFORCE-style \citep{williams1992simple} optimization framework. 
Despite these advances, the application of RL to LLM-driven search engine interactions and reasoning remains largely unexplored.
